\documentclass[11pt]{article}
\usepackage[margin=.82in]{geometry}
\usepackage[T1]{fontenc}
\usepackage{graphicx,booktabs,amsmath,xurl}
\usepackage[hidelinks]{hyperref}
\title{CS 498 Final Project - Track 2\\Complete 3D Scene Reconstruction}
\author{Your Name / NetID}
\date{Fall 2026}
\begin{document}
\maketitle
\section{Method and motivation}
Explain your changes relative to the provided baseline, relevant related work,
and the fully automatic input-to-output pipeline. Category prediction is not required.
We value original ideas and thoughtful experimentation. State your hypothesis,
design rationale, and own contribution; an informative negative result is useful.
Present one final method, selected by the highest mean benchmark across all
ten Val scenes. Identify its exact code/configuration/model snapshot; use that
same frozen version for the final Val table, Test-results ZIP and Kaggle CSV.
% Use e.g. \cite{xia2026fire3d}; reusable entries are in references.bib.
% Add your own related work, then uncomment the bibliography commands below.
\section{Experimental protocol}
Give code revision, dependencies, hardware, checkpoint sources/hashes, commands,
configurations, inference time, and the use of Dev versus Val. Report the
official Kaggle score with its submission ID, not an invented local Test score.
\section{Quantitative results and controlled comparisons}
Report all ten Val cases, split means, the baseline comparison, and at least
one controlled comparison isolating each claimed improvement. Missing/failed
cases count as zero. Report all six metric groups, including static collision-free recall.
Comparisons support the selected final method; do not present several final
entries or combine the best cases/metrics from different versions.
For the optional novel-idea bonus, explain the substantive idea implemented in
that final version and give its actual Val results and controlled comparison.
The same final Kaggle entry completes the private-Test evidence; do not invent
GT comparisons or submit an extra method.
\section{Analysis and visualization}
\subsection{Make the visual story easy to follow}
Use the handout's opening figures as a presentation example, not results to
copy. Start with a few input RGB frames and the known camera poses. Then show
your instance IDs consistently across projected instance views, the assembled
scene, and zoomed isolated-object renders. Two opposite-side views help reveal
missing backs, undersides, or texture; a front-facing image alone is not proof
of a complete reconstruction. Mask/box panels may help explain your method,
but are not additional required output files.

For Dev/Val collision evidence, mark the reported object IDs and geometric
interpenetration witnesses, or state explicitly that no collision was found.
Include proxy/fallback and timeout status where applicable. This is static
geometry evaluation, not a claim about dynamic stability. State whether each
panel is your prediction, a baseline, or a GT reference. Do not edit geometry
or images to fabricate a clean result.
These layout suggestions do not add new scoring criteria or change the
required nine-scene evidence below; Test remains prediction-only.
\subsection{Nine-scene results and visual evidence}
List three distinct selected scene IDs from each split in this table. All
selected outputs must come from the same final method named above. There is
no baseline-score threshold for this evidence component.
\begin{center}
\begin{tabular}{llll}\toprule
Split & Selected scene 1 & Selected scene 2 & Selected scene 3 \\
\midrule
Dev & & & \\
Val & & & \\
Test & & & \\
\bottomrule\end{tabular}
\end{center}
For each selected Dev/Val scene, show compact matching, scene geometry and
appearance, object geometry and appearance, and static collision evidence.
Link the complete evaluator dump, identify instance IDs, and report all six
metric groups. Honest failures must remain visible.
For each selected Test scene, show prediction-only instance, whole-scene,
and isolated-object views (use \path{scripts/student/preview_outputs.py}). Name the scenes for
the official collision check; do not claim private GT comparisons or local
Test metrics. The Kaggle total may be quoted with its submission ID. No post-evaluation report
resubmission is required.
\subsection{Interpretation and failure analysis}
For Dev/Val, show instance matches, full-scene and individual-object GT/prediction views,
occluded-side errors, and collision/contact evidence. Discuss failure
cases, trade-offs, and limitations; do not cherry-pick only successful cases.
\section{AI-human collaboration}
Summarize the AI tools/tasks, important suggestions, your decisions, verification,
failed attempts, and what you learned. Reference the submitted collaboration
record; you remain responsible for all code and reported results.
\section{Reproduction and submission checklist}
List the exact final command, model/checkpoint file SHA-256, ten-case Val
summary, selected scene IDs, per-case evidence, source snapshot, and AI record.
Submit report, actual source code, configuration/dependencies, model declaration,
AI record and the generated HF pointer in the small course ZIP. Large results
and the exact Kaggle CSV/receipt go in your own HF dataset's \path{results.zip}.
State the dataset ID, full 40-character commit, results ZIP SHA-256 and actual
Kaggle submission ID explicitly below; they must match \path{results_location.json}.
Do not use a moving \texttt{main} link or a short hash. Keep the repo private during
development, then make it public/ungated at final hand-in. Verify anonymous access;
do not submit tokens. Keep this revision available through grading/appeals.
An online code link does not replace the submitted source snapshot. Retain both
archive hashes. Explain failures; do not link only to files on your computer.
\begin{description}
\item[HF dataset:] Your account / results repository.
\item[Full commit:] The 40-character commit containing \path{results.zip}.
\item[Results SHA-256:] The generated 64-character file hash.
\item[Kaggle submission:] The actual selected final submission ID.
\end{description}
% \bibliographystyle{unsrt}
% \bibliography{references}
\end{document}
